(1) 轴段 $dx$ 的极转动惯量为 $\rho J_pdx$，其动能密度 $\rho J_p\theta_t^2/2$；小扭角的弹性能密度为 $GJ_p\theta_x^2/2$。对Lagrange密度变分得 $\boxed{\rho J_p\theta_{tt}=GJ_p\theta_{xx}}$，因此扭转波速 $c_t=\sqrt{G/\rho}$，在均匀轴中与截面大小抵消。

(2) 取 $\theta=A\cos(kx-\omega t)$，色散 $\omega=c_t|k|$。内扭矩 $\mathcal T=GJ_p\theta_x=-GJ_pkA\sin(kx-\omega t)$；横截面角速度 $\theta_t=\omega A\sin(kx-\omega t)$，对于向右行波两者瞬时反号，表示正功率沿正 $x$。

(3) 两端自由意味着 $\mathcal T(0)=\mathcal T(L)=0$，即 $\theta_x=0$。本征形状 $\theta_n=\cos(n\pi x/L)$，频率 $\boxed{\omega_n=c_tn\pi/L}$，$n=0,1,\ldots$。$n=0$ 时角度对空间均匀且不产生扭转应变，是整体自由刚体角度或匀速旋转模，不是弹性回复振动。

(4) 对PDE乘 $\theta_t$，得到 $\partial_t[\rho J_p\theta_t^2/2+GJ_p\theta_x^2/2]+\partial_x[-GJ_p\theta_x\theta_t]=0$。平面波平均能密度为 $\bar e=\rho J_p\omega^2A^2/2$，平均能流为 $\bar S=GJ_pk\omega A^2/2$，两者之比 $Gk/(\rho\omega)=c_t$（正向波）。
